% ===== spotify-home-frames.tex: Spotify's home page, the epsilon-greedy case in three frames, 2026-11-06 =====
% The frames of the case and their talk track. Since 2026-11-09 two decks input this file, each as one section:
% block3-bandits.tex, the deck of its Friday block, and weekend2-cases.tex, the weekend deck.
% Carlos chose the case on 2026-11-06 ("Spotify, Amazon Payments, and LinkedIn Lirank") and asked: "Produce the
% slides, follow the style from the slides for the matrix factorization part". Read as the Spotify deck, the audiobooks
% case after the 10:00 matrix factorization block: (1) the problem, a big number on click 2; (2) the method as a
% diagram whose top line ties it to the block's game, here the mirror of "Once Upon a Mirror" (w2-cx-private,
% 03_bandits: a mirror shows one product each morning, a tap is a click, epsilon-greedy spins a roulette), click 2
% adding a part; (3) the value, an A/B test with group chips and hatched bars. He did not dictate these slides.
% No number is typed here: numbers.tex is written by tools/cases_numbers.py (block "Spotify's home page"), which
% keeps the sentence of each source. The three mixes of frame 2 and their estimates are INVENTED (the paper does not
% list its mixes), and the slide says so. A "Say" comment above each frame is the talk track. Build and check:
%   python3 tools/cases_numbers.py
%   ../style/tools/build.sh block3-bandits.tex
%   python3 tools/check_numbers.py block3-bandits.pdf tools/expected-block3-bandits.txt
% The source line is a link (\weblink, violet) to the free full text whose numbers the slides quote: the arXiv
% version read in full for a paper (the line still names the published venue), the post itself for a blog
% (Carlos, 2026-11-12: "put hyperlinks in the citations of the business cases").
\newcommand{\shsource}{\weblink{https://arxiv.org/abs/2509.05460}{Feijer and others (Spotify), CONSEQUENCES workshop at RecSys 2025}.}

% ===== Frame 1. Music or podcasts =====
% Say: Spotify's home page is a stack of shelves, some of music, some of podcasts. How many of each should a
% listener see? A common rule gives each user the mix of their own past listening. But tastes move: the paper's
% first figure shows the share of podcasts rising and falling over the day and over the week, and the authors write
% that preferences may vary with the time of day, the day of the week, even the device. The rule only looks back.
% [click] And this page is built for more than 600 million users.
\begin{frame}{Spotify's home page: music or podcasts?}
  \begin{tikzpicture}[x=1mm, y=1mm]
    \path[use as bounding box] (0,1) rectangle (130,-54);
    \node[body, anchor=north west, text width=64mm] at (0,0)
      {The home page shows shelves of music and podcasts. A common rule gives users the mix of their past
       listening.\par\vspace{2.4mm}
       But tastes can change with the hour, the day of the week, even the device.};
    \node[big, anchor=north west, visible on=<2->] at (76,-1) {Over \shUsers{} million};
    \node[body, anchor=north west, text width=54mm, visible on=<2->] at (76,-10) {users of the home page};
    \node[src, anchor=north west, text width=130mm] at (0,-50.6) {\shsource};
  \end{tikzpicture}
\end{frame}

% ===== Frame 2. Epsilon-greedy with a network =====
% Say: The bandit, in the words of our mirror game. The arms are mixes: mostly music, half and half, mostly podcasts.
% These three are my invention; the paper does not list its mixes. The reward is engagement, for example a stream:
% one if the listener plays anything from the page, zero if not, the tap of the mirror. The mirror estimated each product by counting taps.
% Here a network estimates the chance of a stream for every mix, from the listener's history, the device, the time
% of day and the day of the week. It learned from two weeks of logs and was checked on a third. The bandit shows the
% mix with the highest estimate. [click] And it spins a roulette, as the mirror did: on 98.5 percent of the pages
% the highest estimate, on 1.5 percent a random mix near the listener's usual one. Spotify also tried a mix drawn
% at random from all of them; it gave fewer successes to learn from and a weaker model. The authors' explanation:
% such pages often stray too far from what people like.
\begin{frame}{Epsilon-greedy on the home page}
  \begin{tikzpicture}[x=1mm, y=1mm]
    \path[use as bounding box] (0,1) rectangle (130,-54);
    \node[body, anchor=north west, text width=130mm] at (0,0)
      {The mirror counted taps. Here a network estimates the chances.};
    % the network and its inputs
    \node[anchor=north, inner sep=0, font=\small, align=center] at (19,-9) {Listening history,\\ device, time of day};
    \node[net] (net) at (19,-25.5) {Network};
    \draw[flow] (19,-16.8) to (net.north);
    \draw[draw=nib, line width=0.7pt, line join=miter] (net.east) to (40,-25.5);
    \draw[draw=nib, line width=0.7pt, line join=miter] (40,-18) to (40,-33);
    \foreach \my in {-18,-25.5,-33} {\draw[flow] (40,\my) to (44.4,\my);}
    % the estimates, INVENTED (tools/cases_numbers.py, SH_EST): one bar per mix, the highest hatched
    \node[measure, anchor=north west] at (46,-9) {Chance of a stream from the page};
    \foreach \my/\mname in {-18/Mostly music, -25.5/Half and half, -33/Mostly podcasts}
      {\node[body, anchor=west] at (46,\my) {\mname};}
    \draw[bar plain] (78,-20.5) rectangle ({78+\shEstMusic},-15.5);
    \fill[hatch] (78,-28) rectangle ({78+\shEstHalf},-23);
    \draw[brand, line width=0.7pt, line join=miter] (78,-28) rectangle ({78+\shEstHalf},-23);
    \draw[bar plain] (78,-35.5) rectangle ({78+\shEstPod},-30.5);
    % the roulette (click 2)
    \node[body, anchor=north west, text width=130mm, visible on=<2->] at (0,-38.6)
      {\textbf{\shGreedy\%} of the pages: the highest estimate.};
    \node[body, anchor=north west, text width=130mm, visible on=<2->] at (0,-43.4)
      {\textbf{\shExplore\%}: a random mix near the usual one.};
    \node[src, anchor=north west, text width=130mm] at (0,-50.6) {Invented mixes and estimates. \shsource};
  \end{tikzpicture}
\end{frame}

% ===== Frame 3. The A/B test =====
% Say: The value. Spotify ran an A/B test: the bandit against the mix of the last 90 days, which had beaten a
% 7-day version of the same rule offline. [click] The measure: streams per podcast shelf shown. It rose by 36.6
% percent, a relative change. Over shelves of both kinds the rise was 3.93 percent. [click] And people listened
% more: minutes listened on the whole platform rose by 1.28 percent, and the share of users who started a stream
% from the home page by 1.54 percent. The bandit was in production as of March 2025. The paper gives neither
% the size nor the length of the test. Among the authors' plans: to try UCB, which "has been shown to
% outperform epsilon-greedy strategies in many settings".
\begin{frame}{The value: an A/B test}
  \begin{tikzpicture}[x=1mm, y=1mm]
    \path[use as bounding box] (0,1) rectangle (130,-54);
    \node[chip, font=\small, text width=26mm, align=center, minimum width=28mm, minimum height=7.6mm, inner sep=1.5pt]
      (users) at (14.5,-24) {Home page\\ users};
    \node[lbl, anchor=north] at (14.5,-29.4) {split at random};
    \foreach \gy/\gname/\gtext in {-15/g1/{The mix of the last \shDays{} days}, -33/g2/{The bandit}}
      {\node[chip, font=\small, text width=33mm, align=center, minimum width=36mm, minimum height=7.6mm,
         inner sep=1.5pt] (\gname) at (52,\gy) {\gtext};}
    \draw[draw=nib, line width=0.7pt, line join=miter] (users.east) to (30,-24);
    \draw[draw=nib, line width=0.7pt, line join=miter] (30,-15) to (30,-33);
    \foreach \gy in {-15,-33} {\draw[flow] (30,\gy) to (33.9,\gy);}
    % the podcast measure and its bars (click 2): the control's bar is 22 mm long
    \node[measure, anchor=north west, text width=56mm, visible on=<2->] at (74,-3)
      {Streams per podcast shelf shown, relative change};
    \draw[bar plain, visible on=<2->] (74,-17.5) rectangle (96,-12.5);
    \fill[hatch, visible on=<2->] (74,-35.5) rectangle ({74+22*\shIdxPod/100},-30.5);
    \draw[brand, line width=0.7pt, line join=miter, visible on=<2->] (74,-35.5) rectangle ({74+22*\shIdxPod/100},-30.5);
    \node[big, anchor=west, visible on=<2->] at ({76.5+22*\shIdxPod/100},-33) {+\shPod\%};
    % minutes listened and production (click 3)
    \node[anchor=north west, inner sep=0, font=\large, visible on=<3->] at (74,-40) {+\shMinutes\%};
    \node[measure, anchor=north west, text width=37mm, visible on=<3->] at (93,-40.6)
      {minutes listened on Spotify, relative change};
    \node[body, anchor=north west, text width=64mm, visible on=<3->] at (0,-40) {In production as of \shDeployed.};
    \node[src, anchor=north west, text width=130mm] at (0,-50.6) {\shsource};
  \end{tikzpicture}
\end{frame}
